\documentclass[10pt,a4paper]{amsart}
\usepackage[margin=19mm]{geometry}
\usepackage{amsmath,amssymb,mathtools}
\usepackage[scr=rsfs,cal=euler]{mathalfa}
\usepackage{xfrac}
\usepackage[unicode,hidelinks,bookmarks=false]{hyperref}
\newcommand{\ie}{i.e.\ }
\newcommand{\cf}{cf.\ }
\newcommand{\resp}{respectively\ }
\newcommand{\Subsection}{\subsection}
\providecommand{\nobreakdash}{\nobreak\hbox{-}}
\setlength{\emergencystretch}{2em}
\multlinegap0pt
\usepackage{bm}
\usepackage{bbm}
\usepackage{dsfont}
\usepackage{xspace}
\usepackage{cancel}
\usepackage{mathtools}
\usepackage{amsthm}
\makeatletter
\@ifundefined{theorem}{\newtheorem{theorem}{Theorem}}{}
\@ifundefined{corollary}{\newtheorem{corollary}[theorem]{Corollary}}{}
\@ifundefined{lemma}{\newtheorem{lemma}[theorem]{Lemma}}{}
\@ifundefined{proposition}{\newtheorem{proposition}[theorem]{Proposition}}{}
\makeatother
\title[On the smallest numerical semigroups closed under affine map]{On the smallest numerical semigroups closed under affine maps}
\author{Amelia \'Alvarez, Carlos-Jes\'us Moreno-\'Avila, Ignacio Ojeda}
\date{Source: arXiv:2607.03258v1 (CC BY 4.0)}
\begin{document}
\maketitle

\noindent\emph{Source and licence.} On the smallest numerical semigroups closed under affine maps. Version arXiv:2607.03258v1 (CC BY 4.0), distributed under the Creative Commons Attribution 4.0 International licence, as recorded in the arXiv licence metadata for this exact version and shown on the abstract page https://arxiv.org/abs/2607.03258v1. Full rights notice: licenses/arxiv-2607.03258-CC-BY-4.0.txt.\par\smallskip

\noindent\textit{Editorial scope.} Counted unit: the Theorem whose source label is \texttt{th:arith\_family\_PF} in the article (source0.tex lines 1344--1360, source label \texttt{th:arith\_family\_PF}), delivered together with the article's own complete original proof of it (source0.tex lines 1362--1426, one \texttt{proof} environment of 65 lines). No mathematics is written, repaired or re-derived: statement and proof are the article's own text line for line. Required context retained uncounted: prop:PF1 (lines 728--735); th:maxcrit\_mA (lines 822--824); cor:mminus1\_maximal (lines 832--834); cor:two\_layers (lines 880--885); cor:two\_bands (lines 902--911); lem:splitting-lambda (lines 1221--1226); lem:full-level-maximals (lines 1260--1262). Every definition, display and label of those lines is retained unchanged. Omitted from this package and not redistributed: 1--727, 736--821, 825--831, 835--879, 886--901, 912--1220, 1227--1259, 1263--1343, 1361--1361, 1427--1617. Nothing inside a retained excerpt is omitted or altered: every retained body is delivered exactly as the article's lines are written, so no omission receipt of any kind accompanies this package. Citations: the retained excerpts carry no citation command at all, so no bibliography entry of the article is transcribed and no reference is redistributed. Reference adaptation: none was needed and none was made; every reference inside the delivered unit resolves inside it, so no unresolved reference or citation is produced. Printed slips kept literal: no printed word, symbol, hypothesis or number of the counted statement or of its proof was altered. Layout and class: the article's own class is \texttt{amsart} with packages including amsmath, amssymb, geometry, hyperref, xcolor; the wrapper is the lane's pinned \texttt{amsart} editorial wrapper at 10pt on A4, which restates the theorem-like environments the retained text needs and adds the article's own macro definitions that the retained text needs (none), with extra wrapper packages (bm, bbm, dsfont, xspace, cancel, mathtools). The delivered numbering is the unit's own. Assets: no retained span contains a figure environment, a graphics-inclusion command, a PDF-page inclusion, a table or a TikZ picture; no image file is copied into this package, the package declares no asset dependency and no image file is redistributed with it. Licence: version arXiv:2607.03258v1 of this article is distributed under the Creative Commons Attribution 4.0 International licence, as recorded in the arXiv licence metadata for this exact version and shown on the abstract page https://arxiv.org/abs/2607.03258v1. A full rights notice is delivered at licenses/arxiv-2607.03258-CC-BY-4.0.txt. Characters: every retained excerpt is pure ASCII, so the delivered engine, XeLaTeX with its default Latin Modern text font, reports no missing character.
\begin{proposition}\label{prop:PF1}
Let $S$ be a numerical semigroup and let $m=\operatorname{m}(S)$. Then
\[
\operatorname{PF}(S)=
\{\,w-m \mid w\in \operatorname{Maximals}_{\preceq_S}\operatorname{Ap}(S)\,\},
\]
where $x\preceq_S y$ if and only if $y-x\in S$.
\end{proposition}

\begin{theorem}\label{th:maxcrit_mA}
For $\ell\in\{0,\ldots,m-1\}$, we have $W(\ell) \in \operatorname{Maximals}_{\preceq_S}(\operatorname{Ap}(S))$ if and only if $\mathsf{m}_A(\ell+A_j)\le \mathsf{m}_A(\ell)$ for every $j\in\{1,\ldots,n-1\}$ such that $\ell+A_j<m$.
\end{theorem}

\begin{corollary}\label{cor:mminus1_maximal}
We have $W(m-1) \in \operatorname{Maximals}_{\preceq_S}(\operatorname{Ap}(S))$. In particular, $W(m-1)-m\in\operatorname{PF}(S)$.
\end{corollary}

\begin{corollary}\label{cor:two_layers}
If $W(\ell) \in \operatorname{Maximals}_{\preceq_S}(\operatorname{Ap}(S))$, then
$u_{n-1}^{(\ell)}\in\{u_{n-1}^{(m-1)}-1,u_{n-1}^{(m-1)}\}$.
Equivalently, if $u_{n-1}^{(\ell)}\le u_{n-1}^{(m-1)}-2$, then
$W(\ell)$ is not maximal.
\end{corollary}

\begin{corollary}\label{cor:two_bands}
For $\ell\in\{0,\ldots,m-1\}$, if $W(\ell) \in \operatorname{Maximals}_{\preceq_S}(\operatorname{Ap}(S))$, then exactly one of the following
conditions holds:
\begin{enumerate}
\item[\rm(I)] $u_{n-1}^{(\ell)}=u_{n-1}^{(m-1)}$ and
$r(\ell)\le r(m-1)$;
\item[\rm(II)] $u_{n-1}^{(\ell)}=u_{n-1}^{(m-1)}-1$ and
$r(\ell)\ge r(m-1)+1$.
\end{enumerate}
\end{corollary}

\begin{lemma}\label{lem:splitting-lambda}
Let $k\ge 2$. If $0\le q\le a$, $0\le r<A_k$, and $qA_k+r<A_{k+1}$, then
\[
\mathsf m_{A^{(k)}}(qA_k+r)=\mathsf m_{A^{(k-1)}}(r)+q.
\]
\end{lemma}

\begin{lemma}\label{lem:full-level-maximals}
For every $k\ge 2$, $\mathcal M_{k-1}(A_k)=\{A_k-1-q\mid q=0,\ldots,k-2\}$. Moreover, $\mathsf m_{A^{(k-1)}}(A_k-1-q)=a+q(a-1)$ for $q=0,\ldots,k-2$.
\end{lemma}

\begin{theorem}\label{th:arith_family_PF}
Assume that $m=cA_{n-1}+1$ for some $c\in\{1,\dots,a\}$, and set
$\alpha:=W(m-1)-m=d(m-1)+(c-1)m$. Then
\[
\operatorname{Maximals}_{\preceq_S}\operatorname{Ap}(S)
=
\{\,W(m-1-q)\mid q=0,\ldots,n-2\,\},
\]
and
\[
\operatorname{PF}(S)
=
\{\alpha,\alpha-b,\alpha-2b,\ldots,\alpha-(n-2)b\}.
\]
In particular, $\operatorname{PF}(S)$ is an arithmetic progression of
length $n-1$ with common difference $-b$, and $\operatorname{t}(S)=n-1$.
\end{theorem}

\begin{proof}
If $n=2$, then $A=\{A_1\}=\{1\}$ and
$W(i)=i(d+m)$ for $0\le i\le m-1$. Hence
$W(i)\preceq_S W(j)$ whenever $0\le i\le j\le m-1$, so
$W(m-1)$ is the unique maximal element of $\operatorname{Ap}(S)$. The conclusion follows from Proposition~\ref{prop:PF1}. Hence assume $n\ge 3$.

We first determine the maximal elements of the Ap\'ery poset. Let
$W(\ell)$ be maximal. Since $m-1=cA_{n-1}$, we have $r(m-1)=0$. By
Corollaries~\ref{cor:two_layers} and~\ref{cor:two_bands}, either
$\ell=m-1$, or $\ell=(c-1)A_{n-1}+r$ with $1\le r\le A_{n-1}-1$.

In the latter case, since $
\ell+A_{n-1}=cA_{n-1}+r=(m-1)+r\ge m$,
the index $j=n-1$ is not admissible in Theorem~\ref{th:maxcrit_mA}.
Thus, for every $j\in\{1,\ldots,n-2\}$ such that $r+A_j<A_{n-1}$, we have
$
\ell+A_j=(c-1)A_{n-1}+(r+A_j)<cA_{n-1}+1=m$,
and hence Theorem~\ref{th:maxcrit_mA} gives
$
\mathsf m_A(\ell+A_j)\le \mathsf m_A(\ell).
$
Applying Lemma~\ref{lem:splitting-lambda} to both sides, we obtain $
\mathsf m_{A^{(n-2)}}(r+A_j)\le \mathsf m_{A^{(n-2)}}(r)$.
Therefore $r\in \mathcal M_{n-2}(A_{n-1})$. Hence, by
Lemma~\ref{lem:full-level-maximals}, $r=A_{n-1}-1-p$ for some
$p=0,\ldots,n-3$. Thus $\ell=cA_{n-1}-1-p=m-1-(p+1)$.

Conversely, $W(m-1)$ is maximal by
Corollary~\ref{cor:mminus1_maximal}. Let $q\in\{1,\ldots,n-2\}$, and set
$\ell_q:=m-1-q=(c-1)A_{n-1}+r_q$, where $r_q=A_{n-1}-q$. By
Lemma~\ref{lem:full-level-maximals}, $r_q\in\mathcal M_{n-2}(A_{n-1})$.

We verify the criterion in Theorem~\ref{th:maxcrit_mA}. Let
$j\in\{1,\ldots,n-1\}$ and assume $\ell_q+A_j<m$. Then $j\ne n-1$,
because $\ell_q+A_{n-1}=m-1-q+A_{n-1}\ge m$, as
$q\le n-2<A_{n-1}$. Hence $j\le n-2$.

If $r_q+A_j<A_{n-1}$, then
$\mathsf m_{A^{(n-2)}}(r_q+A_j)\le\mathsf m_{A^{(n-2)}}(r_q)$, and
Lemma~\ref{lem:splitting-lambda} gives
$\mathsf m_A(\ell_q+A_j)\le\mathsf m_A(\ell_q)$. If
$r_q+A_j=A_{n-1}$, then $\ell_q+A_j=m-1$, and
$\mathsf m_A(m-1)=c\le c+q(a-1)=\mathsf m_A(\ell_q)$, where the last
equality follows from Lemmas~\ref{lem:splitting-lambda} and
\ref{lem:full-level-maximals}. Thus Theorem~\ref{th:maxcrit_mA} shows
that $W(\ell_q)=W(m-1-q)$ is maximal.

Therefore $\operatorname{Maximals}_{\preceq_S}\operatorname{Ap}(S)
=
\{\,W(m-1-q)\mid q=0,\ldots,n-2\,\}$.
By Proposition~\ref{prop:PF1},
$\operatorname{PF}(S)=\{\,W(m-1-q)-m\mid q=0,\ldots,n-2\,\}$.

It remains to identify this set explicitly. By
Lemmas~\ref{lem:splitting-lambda} and~\ref{lem:full-level-maximals},
$\mathsf m_A(m-1-q)=c+q(a-1)$ for $q=0,\ldots,n-2$. Hence
$W(m-2-q)-W(m-1-q)=-d+(a-1)m=-b$ for $q=0,\ldots,n-3$, because
$d=(a-1)m+b$.

Finally, $\mathsf m_A(m-1)=c$, so
$W(m-1)-m=d(m-1)+(c-1)m=\alpha$. Thus
$\operatorname{PF}(S)=
\{\alpha,\alpha-b,\alpha-2b,\ldots,\alpha-(n-2)b\}$, and the result
follows.
\end{proof}

\end{document}
